\chapter{Clean Four-Loss Focal Sweep}

\section{Purpose of the Clean Rerun}

The \texttt{0401-focal-pib-sweep-clean-4loss-cn2-5e14} campaign reran the focal-plane
loss family after the PIB bug and the data-generation cleanup. The central
question was no longer ``can PIB be made large?'' It was ``does a large
focal-plane bucket ratio actually increase detector utility?''

\section{The First-Look Result}

At first glance, the answer looked positive. Several strategies pushed
$\PIB@10\,\mu$m upward:
\runfigure{0401-focal-pib-sweep-clean-4loss-cn2-5e14/10_figA_all_strategies_pib_comparison.png}
  {0401 PIB comparison}
  {All four corrected focal-plane losses compared in terms of
  $\PIB@10\,\mu$m.}

\runfigure{0401-focal-pib-sweep-clean-4loss-cn2-5e14/12_figC_focal_all_strategies.png}
  {0401 focal-plane montage}
  {Focal-plane visual comparison across the four corrected strategies.}

\runfigure{0401-focal-pib-sweep-clean-4loss-cn2-5e14/13_figD_wfrms_all_strategies.png}
  {0401 wavefront summary}
  {Wavefront summary shown only as supporting diagnostics; it is not the
  detector-side objective.}

\section{Detector Utility Collapse}

The decisive comparison is the absolute received power, not the normalized PIB.
\runfigure{0401-focal-pib-sweep-clean-4loss-cn2-5e14/16_throughput_analysis.png}
  {0401 throughput analysis}
  {Throughput collapse across the conventional focal-plane losses.}

\runfigure{0401-focal-pib-sweep-clean-4loss-cn2-5e14/17_received_power_analysis.png}
  {0401 received-power analysis}
  {Detector-side received power reveals that all four conventional strategies
  underperform the turbulent baseline at 10\,$\mu$m.}

\runfigure{0401-focal-pib-sweep-clean-4loss-cn2-5e14/16_absolute_bucket_power_histogram.png}
  {0401 absolute bucket histogram}
  {Absolute bucket power rather than normalized bucket fraction is the key
  detector-side quantity.}

\runfigure{0401-focal-pib-sweep-clean-4loss-cn2-5e14/20_encircled_energy_per_strategy.png}
  {0401 encircled energy}
  {Encircled-energy comparison across the clean four-loss sweep.}

\begin{table}[H]
  \centering
  \begin{tabular}{lccccc}
    \toprule
    \smalltablehead{Strategy} & \smalltablehead{$\CO_{\mathrm{out}}$} & \smalltablehead{$\PIB@10\,\mu$m} & \smalltablehead{$\TP$} & \smalltablehead{$\RP@10\,\mu$m} & \smalltablehead{vs.\ turb.} \\
    \midrule
    Turbulent baseline & -- & 0.8084 & 0.9199 & 0.7443 & -- \\
    PIB only & -- & 0.9116 & 0.5016 & 0.4355 & $-41.5\%$ \\
    Strehl only & -- & 0.8185 & 0.5710 & 0.4449 & $-40.2\%$ \\
    Intensity overlap & -- & 0.2563 & 0.0803 & 0.0199 & $-97.3\%$ \\
    CO+PIB hybrid & 0.9726 & 0.9582 & 0.0165 & 0.0153 & $-97.9\%$ \\
    \bottomrule
  \end{tabular}
  \caption{The clean rerun reverses the naive PIB story. Every conventional
  focal-plane loss underperforms the turbulent baseline in received power at
  10\,$\mu$m.}
\end{table}

\section{Failure Modes}

The strategies fail in different ways. \texttt{PIB only} and \texttt{Strehl only} improve the
bucket ratio but throw away too much energy. \texttt{Intensity overlap} collapses both
throughput and detector utility. \texttt{CO+PIB hybrid} is the most deceptive: it
produces excellent-looking $\CO$ and PIB numbers while almost eliminating the
useful detector signal.

\begin{discovery}
The clean rerun established the main detector-side lesson of the entire report:
normalized sharpness objectives are not detector objectives. A static D2NN can
make the surviving energy look sharp while simultaneously destroying the amount
of energy that reaches the detector.
\end{discovery}

The next step was therefore not to abandon focal-plane optimization, but to
design losses in which throughput is structurally inside the objective.
